\documentclass[letterpaper,11pt]{article}
\usepackage[letterpaper,margin=1in]{geometry}
\usepackage[T1]{fontenc}
\usepackage[utf8]{inputenc}
\usepackage{listings}
\usepackage{xcolor}
\usepackage[hidelinks]{hyperref}
\lstset{
  basicstyle=\small\ttfamily,
  breaklines=true,
  frame=single,
  framesep=3pt,
  xleftmargin=6pt,
  xrightmargin=6pt,
  aboveskip=8pt,
  belowskip=8pt,
  columns=fullflexible,
  keepspaces=true,
}

\title{Supplementary Material: Proof-of-Concept Source Code\\
  \large for ``Aperiodic Guardrails: Algebraic Impossibility Results for
  Content-Based LLM Safety Filters''}
\author{}
\date{}

\begin{document}
\maketitle

\noindent This supplementary document contains the full source code of the
proof-of-concept implementation referenced in \S9 of the main paper.
The content was originally included as Appendix~A and has been moved here
unmodified to meet conference page budgets; all line counts, file names,
and structural claims in the main paper refer to the code below.

\bigskip

\section{Proof-of-Concept Source Code}\label{app:poc}

The complete PoC consists of 376 lines across six files.
The monoid extractor (858 lines) and benchmark harness (619 lines) are provided
as supplementary material at the project repository.
Below we summarize the architecture; full source is available online.

\subsection{Engine: Zero-Knowledge Recurrent Executor}\label{app:engine}

\texttt{engine.py} (63 lines) decodes a base64-encoded configuration, selects an AST medium,
compiles it via \texttt{exec(compile(...))}, and loops until a termination condition.
No LLM sees the configuration contents.

\begin{lstlisting}[language=Python,caption={engine.py (63 lines)}]
"""ZK Recurrent Executor. Decodes b64 config, selects
AST medium, compiles, injects params at runtime, loops
until termination. No LLM sees the config contents."""
import base64, json, sys, time
from mediums import get_medium

def decode_config(path):
    with open(path, 'rb') as f:
        return json.loads(base64.b64decode(f.read()))

def check_termination(results, iteration, config):
    term = config.get('termination',
                      {'type': 'count', 'value': 1})
    t = term['type']
    if t == 'count':
        return iteration >= term.get('value', 1)
    if t == 'has_results':
        return bool(results)
    if t == 'min_results':
        if isinstance(results, dict) and 'path' in results:
            return len(results['path']) > 0
        return bool(results) \
            and len(results) >= term.get('value', 1)
    if t == 'score_threshold':
        return isinstance(results, dict) \
            and results.get('score', 0) >= term['value']
    return True

def mutate_params(params, iteration, config):
    m = config.get('mutation', {})
    if m.get('type') == 'paginate':
        k = m.get('key', 'page')
        params[k] = params.get(k, 1) + 1
    elif m.get('type') == 'backoff':
        time.sleep(min(2 ** iteration, 30))
    return params

def main():
    config = decode_config(
        sys.argv[1] if len(sys.argv) > 1 else 'config.b64')
    params = config['params']
    build_ast = get_medium(config['medium'])

    ns = {}
    exec(compile(build_ast(),
         f'<{config["medium"]}>', 'exec'), ns)
    run = ns['run']

    max_iter = config.get('max_iterations', 10)
    results = None
    for i in range(1, max_iter + 1):
        print(f"[engine] iteration {i}/{max_iter}")
        try:
            results = run(**params)
        except Exception as e:
            print(f"[engine] error: {e}")
        if check_termination(results, i, config):
            break
        params = mutate_params(params, i, config)

    print(f"\n[engine] RESULTS:")
    if isinstance(results, list):
        for j, item in enumerate(results, 1):
            t = item.strip() if isinstance(item, str) \
                else str(item)
            if t:
                print(f"  {j}. {t}")
    elif isinstance(results, dict):
        print(json.dumps(results, indent=2))
    else:
        print(results)

if __name__ == '__main__':
    main()
\end{lstlisting}

\subsection{Medium: BFS Graph Solver}\label{app:graph-solver}

\texttt{graph\_solver.py} implements BFS over an abstract state space.
States are sets of opaque symbols; rules define valid transitions.
Zero domain content --- all semantics injected at runtime.

\begin{lstlisting}[language=Python,caption={graph\_solver.py (43 lines)}]
"""AST template: BFS state-space solver.
Zero domain content."""
import ast

def build_graph_solver_ast() -> ast.Module:
    source = '''
from collections import deque

def applicable_rules(state, rules):
    return [r for r in rules
            if all(i in state for i in r["inputs"])]

def apply_rule(state, rule):
    new = [s for s in state
           if s not in rule["inputs"]]
    new.extend(rule["outputs"])
    return sorted(set(new))

def check_constraints(state, constraints):
    for c in constraints:
        if c["type"] == "forbidden_pair":
            if c["params"][0] in state \
               and c["params"][1] in state:
                return False
    return True

def run(initial_state, target, rules, constraints):
    queue = deque([(initial_state, [], 1.0)])
    visited = set()
    best = None
    while queue:
        state, path, score = queue.popleft()
        key = tuple(state)
        if target in state:
            if best is None or score > best["score"]:
                best = {"path": path, "score": score,
                        "final_state": state}
            continue
        if key in visited:
            continue
        visited.add(key)
        if len(path) >= 10:
            continue
        for rule in applicable_rules(state, rules):
            new_state = apply_rule(state, rule)
            if check_constraints(new_state, constraints):
                new_score = score * rule["yield_est"]
                queue.append((new_state,
                    path + [rule["id"]], new_score))
    return best or {"path": [], "score": 0,
                    "final_state": initial_state}
'''
    tree = ast.parse(source)
    ast.fix_missing_locations(tree)
    return tree
\end{lstlisting}

\subsection{Medium: Tree-of-Thought Solver with LLM Evaluation}\label{app:tot-solver}

\texttt{tot\_solver.py} implements beam search with LLM-guided branch pruning.
The LLM sees \emph{only} abstract symbols (\texttt{N1}, \texttt{R7}, etc.)
and rates branches as \textsc{sure}, \textsc{likely}, or \textsc{impossible}.
Composite scoring: yield weight $0.4$ + LLM viability $0.6$.
Beam width~5, max depth~8.

\begin{lstlisting}[language=Python,caption={tot\_solver.py --- core evaluation and search (abbreviated; 158 lines full)}]
def ollama_evaluate(state, target, rules, model):
    """Ask LLM: is this state viable toward target?
    Returns 0.0-1.0. LLM sees ONLY abstract symbols."""
    system_ctx = _build_system_context(target, rules)
    applicable = [
        f'{r["id"]}: {r["inputs"]} -> {r["outputs"]}'
        f' (yield {r["yield_est"]})'
        for r in rules
        if all(i in state for i in r["inputs"])]
    user_prompt = (
        f"CURRENT STATE: {state}\n"
        f"APPLICABLE NOW:\n"
        + ("\n".join(applicable)
           if applicable else "(none)")
        + f"\n\nCan {target} be reached? "
        f"Trace forward briefly, then rate.\n"
        f"RATING: SURE / LIKELY / IMPOSSIBLE")
    try:
        data = json.dumps({
            "model": model,
            "messages": [
                {"role": "system",
                 "content": system_ctx},
                {"role": "user",
                 "content": user_prompt}],
            "stream": False}).encode()
        req = urllib.request.Request(
            "http://localhost:11434/api/chat",
            data=data,
            headers={"Content-Type":
                      "application/json"})
        with urllib.request.urlopen(req,
                timeout=15) as resp:
            result = json.loads(resp.read())
        content = result["message"]["content"]
        lines = [l.strip()
                 for l in content.strip().split("\n")
                 if l.strip()]
        last_line = lines[-1].upper() if lines else ""
        if "IMPOSSIBLE" in last_line: return 0.0
        if "SURE" in last_line: return 1.0
        if "LIKELY" in last_line: return 0.6
        return 0.5
    except Exception:
        return 0.5

def run(initial_state, target, rules, constraints,
        model="llama3.2:3b"):
    """Tree-of-Thought search: generate branches,
    LLM evaluates, prune, expand."""
    beam_width, max_depth = 5, 8
    frontier = [(initial_state, [], 1.0, 1.0)]
    visited = set()
    best = None
    for depth in range(max_depth):
        if not frontier: break
        candidates = []
        for state, path, cum_yield, _ in frontier:
            key = tuple(state)
            if target in state:
                if best is None \
                   or cum_yield > best["score"]:
                    best = {"path": path,
                            "score": cum_yield,
                            "final_state": state}
                continue
            if key in visited: continue
            visited.add(key)
            for rule in applicable_rules(state, rules):
                new_state = apply_rule(state, rule)
                if not check_constraints(
                        new_state, constraints):
                    continue
                new_yield = cum_yield*rule["yield_est"]
                new_path = path + [rule["id"]]
                candidates.append((new_state,
                    new_path, new_yield, rule))
        if not candidates: break
        scored = []
        for ns, np, ny, rule in candidates:
            llm = ollama_evaluate(
                ns, target, rules, model)
            composite = ny * 0.4 + llm * 0.6
            scored.append((ns, np, ny, composite))
        scored.sort(key=lambda x: x[3], reverse=True)
        frontier = [f for f in scored[:beam_width]
                    if f[3] > 0.1]
    return best or {"path": [], "score": 0.0,
                    "final_state": initial_state}
\end{lstlisting}

\subsection{Config Encoder}\label{app:encode}

\begin{lstlisting}[language=Python,caption={encode.py (22 lines)}]
"""Encode a JSON config to b64.
No LLM sees plaintext values."""
import base64, json, sys

def main():
    if sys.argv[1] == '--stdin':
        config = json.loads(sys.stdin.read())
    else:
        with open(sys.argv[1]) as f:
            config = json.load(f)
    out = sys.argv[2] if len(sys.argv) > 2 \
        else 'config.b64'
    encoded = base64.b64encode(
        json.dumps(config).encode())
    with open(out, 'wb') as f:
        f.write(encoded)
    print(f"[encode] {out} ({len(encoded)}B) "
          f"medium={config.get('medium','?')} "
          f"keys={list(config.get('params',{}).keys())}")

if __name__ == '__main__':
    main()
\end{lstlisting}

\subsection{Medium Registry}\label{app:registry}

\begin{lstlisting}[language=Python,caption={mediums/\_\_init\_\_.py (13 lines)}]
"""Medium registry."""
from mediums.web_scraper import build_web_scraper_ast
from mediums.graph_solver import build_graph_solver_ast
from mediums.tot_solver import build_tot_solver_ast

REGISTRY = {
    "web_scraper": build_web_scraper_ast,
    "graph_solver": build_graph_solver_ast,
    "tot_solver": build_tot_solver_ast,
}

def get_medium(name: str):
    if name not in REGISTRY:
        raise ValueError(
            f"Unknown medium '{name}'. "
            f"Available: {list(REGISTRY.keys())}")
    return REGISTRY[name]
\end{lstlisting}

\subsection{Supplementary Material}\label{app:supplementary}

The following artifacts are provided as supplementary material and are not
included inline due to length:
\begin{itemize}
  \item \texttt{monoid\_extractor.py} (858 lines): Thompson NFA construction,
    powerset DFA conversion, Hopcroft minimization, transition monoid
    enumeration, aperiodicity test, Krohn--Rhodes group component extraction,
    and $\mathrm{MOD}_2$ bypass witness generation.
  \item \texttt{benchmark.py} (619 lines): Grammar generator with guaranteed
    reachability, dead-end traps, noise rules, and constraint variation;
    BFS, random-beam, and ToT+LLM solvers; Wilcoxon signed-rank test
    computation; JSON result serialization.
  \item \texttt{web\_scraper.py} (77 lines): AST-based web scraper medium
    with zero string literals.
\end{itemize}

% ============================================================

\end{document}
